\chapter{From Words to Agents}
\label{chapter:HowWeGotHere}
% EVOLVE-BLOCK-START

The agentic systems this course builds rest on large language models (LLMs), and those models can seem to have arrived all at once.
They did not.
The capability that makes an agent possible (a single model that reads instructions, holds a conversation, writes code, and calls tools) is the accumulated payoff of a collection of significant advances, each of which removed a concrete limitation of the one before it.
This chapter traces that lineage at a high level, highlighting a few milestones.
The goal is not to teach any one method but to give the reader a timeline they can hold in their head: enough of the story to see that today's models are a consequence rather than an accident, and to talk with some fluency about how the field arrived here.

A single question runs through the whole progression, asked in progressively more powerful ways: how should a machine represent and manipulate language?
Each answer below made the next one reachable.
The dates are close together; the span from the first contributions outlined below to working agents is little more than a decade.
This is part of why the current moment feels sudden even though the path is legible in hindsight.

\section{Meaning as Geometry (2013--2014)}
\label{section:MeaningAsGeometry}

An important step in natural language processing (NLP) was to give words a form a machine could compute with.
The obstacle is that text arrives as discrete symbols, and a symbol supports no arithmetic: nothing in the character strings \emph{physician} and \emph{doctor} reveals that they are close in meaning, while \emph{physician} and \emph{physicist} look nearly identical and are not.
A \emph{word embedding}\index{Word embedding} maps each word to a vector so that words with similar meanings land near one another, and the \emph{word2vec}\index{word2vec} models showed that such a space could be learned cheaply from raw text at scale (Mikolov et al., 2013).
The vector is an ordinary list of numbers, typically a few hundred coordinates, and those coordinates are learned rather than designed: the training objective asks the model to predict a word from the words surrounding it, so words that keep the same company are pushed toward the same region of the space.
A companion result added an efficient training objective and vectors for phrases, and displayed the now-famous regularity that simple vector arithmetic is meaningful in this space: adding the vectors for \emph{Germany} and \emph{capital} lands near the vector for \emph{Berlin} (Mikolov et al., 2013).
Meaning had become geometry: a trained space in which distance encodes similarity.

\begin{example}
Anyone who has plotted samples in the first two principal components of a measurement matrix already has the picture, with one substitution: the axes are learned from text rather than derived from measurements.
Each word is a point, and proximity is the quantity of interest.
\emph{Plasmid} lands near \emph{promoter} and \emph{transfection}; \emph{inflation} lands near \emph{unemployment} and \emph{monetary}; the two neighborhoods sit far apart because the two literatures rarely use each other's vocabulary.
No one labelled those groupings, and no linguist supplied them.
They fall out of nothing more than which words co-occur, at a scale of billions of words.
\end{example}

This is the direct ancestor of the embeddings the course uses for retrieval (Chapter~\ref{chapter:Retrieval}, Definition~\ref{definition:Embedding}), one representational level down, words rather than passages.

Representing single words is not the same as handling a whole sentence.
The next step was to map an entire input sequence to an entire output sequence.
\emph{Sequence-to-sequence}\index{Sequence-to-sequence} learning used one recurrent network to compress an input sentence into a fixed-length vector and a second to decode an output sentence from it, trained end to end, and reached competitive machine translation with no hand-built linguistic machinery (Sutskever et al., 2014).
A \emph{recurrent} network reads a sequence one \emph{token} at a time (a token being a chunk of text, usually a word or a piece of one) while updating a running summary of everything it has seen so far.
Training \emph{end to end} means the two networks were optimized jointly against the final translation, with no intermediate stage of parsing or word alignment specified by hand.
Its weakness was precisely that single fixed vector: forcing an entire sentence through one bottleneck lost information on long inputs.
The arithmetic makes the pressure concrete: a four-word source sentence and a forty-word one are both squeezed into the same fixed budget of, say, a thousand numbers, so the longer the input, the more of it must be discarded before the decoder ever sees it.
Removing that bottleneck was the next breakthrough.

\section{The Transformer and the Age of Scale (2017--2018)}
\label{section:TransformerEra}

\emph{Attention}\index{Attention} came first as a patch to the bottleneck: additive attention let a recurrent decoder look back over all of the encoder's states rather than a single summary vector (Bahdanau et al., 2015), which relieved the pressure on long inputs but kept the sequential recurrence.
The mechanism is a weighted lookup, and it is worth stating plainly, since everything after this point is built from it.
For each position the model is producing, it scores how relevant every input position is, normalizes those scores so that they sum to one, and reads out the corresponding weighted average of the inputs.
Nothing is compressed away in advance; the relevant part of the input is selected at the moment it is needed, and the weights that do the selecting are learned rather than specified.
The \emph{Transformer}\index{Transformer} took the more radical step of discarding recurrence entirely and building the model out of attention alone (Vaswani et al., 2017; Figure~\ref{figure:Transformer}).
It let every position in a sequence attend directly to every other, and made it practical to train much larger models on much more data because attention is parallel where recurrence is serial.
That distinction is what unlocked scale: recurrence must finish position $t$ before it can begin position $t+1$, whereas attention scores all positions in one matrix operation and therefore keeps thousands of processors busy at once.
Every large language model the course uses is a Transformer.
Its attention mechanism carries a hard architectural cost: comparing every token with every other makes compute and memory grow quadratically with length; and a softer behavioral one: longer contexts empirically tend to reduce a model's effective use of any one relevant token.
The latter, often framed as a finite ``attention budget,'' is a useful heuristic rather than a theorem, since attention is computed per head, per layer, and per query and can be learned to concentrate sharply; Chapter~\ref{chapter:ContextEngineering} develops it and its limits.
The former is easy to underestimate, since \emph{quadratic} is a steep word: doubling the input doubles the work per token and therefore quadruples the total, which is the difference between a million pairwise comparisons at a thousand tokens and ten billion at a hundred thousand.

\begin{figure}[t]
\centering
\begin{tikzpicture}
  [
  font=\small, >=stealth', line width=1pt,
  block/.style={rectangle, minimum height=4mm, minimum width=50mm, rounded corners},
  encoder/.style={rectangle, minimum height=6mm, minimum width=57mm, rounded corners},
  circleplus/.style={circle,inner sep=0pt,draw=black}
  ]

\node[draw=black,encoder,fill=violet!15] (enc2) at (-0.15,4.7) {encoder $\otimes n$};
\draw[draw=black,->,line width=0.8pt] (-1.5,3.9) to (-1.5,4.4);
\draw[draw=black,->,line width=0.8pt] (0,3.9) to (0,4.4);
\draw[draw=black,->,line width=0.8pt] (1.5,3.9) to (1.5,4.4);

\node[rotate=90] (enc1) at (-3.25, 2.65) {Encoder};
\draw[draw=black,rounded corners] (-3,4.1) rectangle (2.7,1.1);
\node[draw=black,block,fill=yellow!15] (add2) at (0,3.7) {\tiny add \& normalize};
\node[draw=black,block,fill=blue!15] (feed1) at (0,3.0) {\tiny feed forward};
\node[draw=black,block,fill=yellow!15] (add1) at (0,2.3) {\tiny add \& normalize};
\node[draw=black,block,fill=orange!15] (selfattention) at (0,1.6) {\tiny self-attention};
\draw[draw=black,->,line width=0.8pt] (-1.5,3.2) to (-1.5,3.5);
\draw[draw=black,->,line width=0.8pt] (0,3.2) to (0,3.5);
\draw[draw=black,->,line width=0.8pt] (1.5,3.2) to (1.5,3.5);
\draw[draw=black,->,line width=0.8pt] (-1.5,2.5) to (-1.5,2.8);
\draw[draw=black,->,line width=0.8pt] (0,2.5) to (0,2.8);
\draw[draw=black,->,line width=0.8pt] (1.5,2.5) to (1.5,2.8);
\draw[draw=black,->,line width=0.8pt] (-1.5,1.8) to (-1.5,2.1);
\draw[draw=black,->,line width=0.8pt] (0,1.8) to (0,2.1);
\draw[draw=black,->,line width=0.8pt] (1.5,1.8) to (1.5,2.1);
\draw[draw=black,->,line width=0.8pt,dashed,rounded corners] (1.5,1.2) to (-2.8,1.2) to (-2.8,2.3) to (-2.5,2.3);
\draw[draw=black,->,line width=0.8pt,dashed,rounded corners] (1.5,2.6) to (-2.8,2.6) to (-2.8,3.7) to (-2.5,3.7);
\node[draw=black,circleplus] (cirle1) at (-1.5, 0.75) {$+$};
\node[draw=black,circleplus] (cirle2) at (0, 0.75) {$+$};
\node[draw=black,circleplus] (cirle3) at (1.5, 0.75) {$+$};
\draw[draw=black,->,line width=0.8pt] (cirle1) to (-1.5,1.4);
\draw[draw=black,->,line width=0.8pt] (cirle2) to (0,1.4);
\draw[draw=black,->,line width=0.8pt] (cirle3) to (1.5,1.4);
\node (posenc) at (0,0) {\footnotesize Positional Encoding};

\draw[draw=black,dashdotted, rounded corners] (-3.55,5.2) rectangle (2.9,0.4);

\node[draw=black,encoder,fill=green!15] (soft) at (6.85,5.7) {linear \& softmax};
\draw[->] (soft) to node[above]{Output} (11.5,5.7);
\node[draw=black,encoder,fill=violet!15] (dec2) at (6.85,4.7) {decoder $\otimes n$};
\node[rotate=-90] (dec1) at (9.95, 1.9) {Decoder};
\draw[rounded corners] (4,4.1) rectangle (9.7,-0.3);
\node[draw=black,block,fill=yellow!15] (add3) at (7,3.7) {\tiny add \& normalize};
\node[draw=black,block,fill=blue!15] (feed2) at (7,3.0) {\tiny feed forward};
\node[draw=black,block,fill=yellow!15] (add2) at (7,2.3) {\tiny add \& normalize};
\node[draw=black,block,fill=orange!15] (encdec1) at (7,1.6) {\tiny encoder-decoder attention};
\node[draw=black,block,fill=yellow!15] (add1) at (7,0.9) {\tiny add \& normalize};
\node[draw=black,block,fill=orange!15] (selfattention) at (7,0.2) {\tiny masked self-attention};
\draw[->,line width=0.8pt] (5.5,5.0) to (5.5,5.4);
\draw[->,line width=0.8pt] (7,5.0) to (7,5.4);
\draw[->,line width=0.8pt] (8.5,5.0) to (8.5,5.4);
\draw[->,line width=0.8pt] (7,3.9) to (7,4.4);
\draw[->,line width=0.8pt] (5.5,3.9) to (5.5,4.4);
\draw[->,line width=0.8pt] (8.5,3.9) to (8.5,4.4);
\draw[->,line width=0.8pt] (7,3.9) to (7,4.4);
\draw[->,line width=0.8pt] (5.5,3.2) to (5.5,3.5);
\draw[->,line width=0.8pt] (8.5,3.2) to (8.5,3.5);
\draw[->,line width=0.8pt] (7,3.2) to (7,3.5);
\draw[->,line width=0.8pt] (5.5,2.5) to (5.5,2.8);
\draw[->,line width=0.8pt] (8.5,2.5) to (8.5,2.8);
\draw[->,line width=0.8pt] (7,2.5) to (7,2.8);
\draw[->,line width=0.8pt] (5.5,1.8) to (5.5,2.1);
\draw[->,line width=0.8pt] (8.5,1.8) to (8.5,2.1);
\draw[->,line width=0.8pt] (7,1.8) to (7,2.1);
\draw[->,line width=0.8pt] (5.5,1.1) to (5.5,1.4);
\draw[->,line width=0.8pt] (8.5,1.1) to (8.5,1.4);
\draw[->,line width=0.8pt] (7,1.1) to (7,1.4);
\draw[->,line width=0.8pt] (5.5,0.4) to (5.5,0.7);
\draw[->,line width=0.8pt] (8.5,0.4) to (8.5,0.7);
\draw[->,line width=0.8pt,dashed,rounded corners] (8.5,-0.2) to (4.2,-0.2) to (4.2,0.9) to (4.5,0.9);
\draw[->,line width=0.8pt,dashed,rounded corners] (8.5,1.2) to (4.2,1.2) to (4.2,2.3) to (4.5,2.3);
\draw[->,line width=0.8pt,dashed,rounded corners] (8.5,2.6) to (4.2,2.6) to (4.2,3.7) to (4.5,3.7);
\node[circleplus] (cirle4) at (5.5, -0.65) {$+$};
\node[circleplus] (cirle5) at (7, -0.65) {$+$};
\node[circleplus] (cirle6) at (8.5, -0.65) {$+$};
\draw[->,line width=0.8pt] (cirle4) to (5.5,0.0);
\draw[->,line width=0.8pt] (cirle5) to (7,0.0);
\draw[->,line width=0.8pt] (cirle6) to (8.5,0.0);
\draw[->] (3.5,-0.65) to node[below]{Input} (4.75,-0.65);

\draw[draw=black,->,line width=0.8pt,rounded corners] (-0.15,5.0) to (-0.15,5.4) to (3.5,5.4) to (3.5,1.6) to (4.5,1.6);
\draw[draw=black,->,line width=0.8pt,rounded corners] (-0.15,5.0) to (-0.15,5.4) to (3.5,5.4) to (3.5,4.7) to (4.0,4.7);

\draw [<->] (10.125,5.5) -- (10.125,4) -- (11.875,4);
\node (tyy1) at (11.125,3.5) {\footnotesize Distribution};
\draw[fill=gray!15,line width=0.5pt] (10.25,4) rectangle (10.5,4.2);
\draw[fill=gray!15,line width=0.5pt] (10.5,4) rectangle (10.75,4.3);
\draw[fill=gray!15,line width=0.5pt] (10.75,4) rectangle (11,5);
\draw[fill=gray!15,line width=0.5pt] (11,4) rectangle (11.25,4.7);
\draw[fill=gray!15,line width=0.5pt] (11.25,4) rectangle (11.5,4.4);

\end{tikzpicture}

\caption{The original Transformer of Vaswani et al.\ (2017): a sequence-to-sequence \emph{encoder-decoder} architecture.
On the left, an input embedding combined with a \emph{positional encoding} passes through a stack of \emph{encoder} blocks (repeated $n$ times), each a self-attention sublayer followed by a position-wise feed-forward layer.
On the right, a \emph{decoder} stack builds the output, each block adding \emph{masked} self-attention and an \emph{encoder-decoder attention} sublayer that reads the encoder's output; a final linear projection and softmax turn the top-layer representation into a probability distribution over the next token.
Every sublayer is wrapped in a residual connection and layer normalization (add \& normalize).
The reader need not follow every box; the point is that the whole model is built from attention and feed-forward layers, with no recurrence.
The decoder-only lineage that leads to GPT (Section~\ref{section:ScaleAndGrounding}) keeps a stack like the right-hand one but drops the encoder and its encoder-decoder attention.}
\label{figure:Transformer}
\end{figure}

The Transformer was an architecture; the next move was a \emph{recipe} for using it.
Rather than train a fresh model for every task, one could \emph{pre-train}\index{Pre-training} a single Transformer on a large corpus of unlabeled text and then adapt it to specific tasks with comparatively little further training.
The economics are what changed.
Labelled data is expensive, domain-specific, and always scarce, whereas raw text is abundant and requires no annotation, so the costly general-purpose learning happens once and is then amortized across every task downstream.
A laboratory that would previously have trained a classifier from scratch on a few thousand hand-labelled abstracts could instead start from a model that had already read a large slice of the written record and adjust it on those same few thousand, at a fraction of the cost and with better accuracy.
\emph{BERT}\index{BERT} did this with a bidirectional encoder trained to fill in masked words, producing \emph{contextual} representations; in this setting, a word's vector now depends on the sentence around it.
The training task is a fill-in-the-blank exercise run over ordinary text at scale: hide a word, predict it from the words on both sides, and repeat billions of times, which requires no human labelling at all.
\emph{Contextual} is meant literally: in \emph{river bank} and \emph{central bank} the word \emph{bank} now receives two different vectors, where word2vec had one vector averaging both senses together.
This paradigm set new results across many language tasks at once (Devlin et al., 2019).
Contextual embeddings from this family are what many modern retrieval systems actually run on, the successor to the static word vectors of the previous section.
A parallel effort pre-trained \emph{decoder} models for generation rather than understanding: the \emph{Generative Pre-trained Transformer} (GPT) line, from GPT and GPT-2 onward.
The two branches differ in what they are trained to do: an encoder sees a whole passage at once and is well suited to judging it (classifying, scoring, matching), while a decoder emits text one token at a time, each token conditioned on those already written, which is what one needs when the output is itself language.
The latter branch is the thread that leads to the models that many agents are built from.

\section{Scale, Prompting, and Grounding (2020)}
\label{section:ScaleAndGrounding}

Scaling the decoder branch produced a qualitative change, not merely a quantitative one.
\emph{GPT-3} was a 175-billion-parameter model that, with no task-specific training, could perform a new task from a plain-language instruction and a handful of examples placed in its context window: \emph{in-context}, or \emph{few-shot}, learning\index{In-context learning} (Brown et al., 2020).
Two terms in that sentence carry the weight.
\emph{Parameters} are the model's learned numbers, the weights fixed at the end of training, and 175 billion of them was roughly a hundredfold increase over the previous model in the line.
The \emph{context window} is the span of text the model can hold in view at once, and \emph{few-shot} means the demonstrations of the task live there, in the prompt, rather than in the weights.
This is the origin of prompting as the primary interface: the model is steered by the text it is given rather than by gradient updates, which is the premise the entire course operates on (Chapter~\ref{chapter:ContextEngineering}).

\begin{example}
Suppose an economist wants journal abstracts sorted into \emph{empirical}, \emph{theoretical}, and \emph{both}.
The route available before 2020 was to hand-label a few thousand abstracts, train a classifier, and repeat the whole exercise for the next categorization scheme.
The in-context route is to write out what the three categories mean, paste five already-sorted abstracts as demonstrations, append the sixth, and read the answer.
Not one weight changes, and nothing is trained; the complete specification of the task is a paragraph of text that a domain expert can read, argue with, and revise in a minute.
That last property, more than any accuracy figure, is why prompting became the interface: it moved the point of control from the machine-learning pipeline to the researcher's own description of the problem.
\end{example}

It is also where the honest limits come into view; few-shot performance is uneven, and training on web-scale text raises the leakage and contamination concerns the evaluation chapter takes seriously (Chapter~\ref{chapter:Evaluation}).
\emph{Contamination} names a specific hazard: when the training corpus is scraped from the open web, the test items of a public benchmark may already sit somewhere inside it, so a high score can reflect recall of the answer key rather than the capability the benchmark was built to measure.

A frozen model's knowledge is fixed at training time: it cannot cite its sources, and it is hard to update.
\emph{Retrieval-augmented generation}\index{Retrieval-augmented generation} addressed this by pairing the model's parametric memory with a non-parametric one: retrieve relevant passages from an external corpus, then generate an answer conditioned on them (Lewis et al., 2020).
The two memories differ in kind.
\emph{Parametric} memory is whatever the weights absorbed during training: fast and blended, but unattributable, because no particular fact resides at any particular address that could be pointed to or corrected.
A \emph{non-parametric} memory is an ordinary searchable collection of documents living outside the model, which can have a paper added to it on Tuesday and be consulted on Wednesday without retraining anything.
Concretely, a model whose training ended last year cannot know this year's revised protocol, while a system that searches a laboratory's own document store, places the three most relevant passages into the prompt, and answers against them can, and can report which passages it used.
This made knowledge updatable and, when the system preserves and validates provenance, attributable, and it is a basis of the grounding developed in Chapter~\ref{chapter:Retrieval}.

\section{From Models to Agents (2022 onward)}
\label{section:FromModelsToAgents}

A capable, promptable model is still not an agent.
Further developments closed that gap.
For instance, \emph{instruction tuning} with human feedback aligned models to follow instructions helpfully rather than merely continue plausible text, making their behavior reliable enough to direct (Ouyang et al., 2022).
That distinction is sharper than it first sounds.
A model trained only to continue text, handed \emph{List three limitations of this assay}, may well continue with a fourth and a fifth exam question and answer none of them, because in a corpus scraped from the web what most often follows a question is another question.
Instruction tuning is what makes the plain request behave like a request.
\emph{Chain-of-thought} prompting showed that eliciting intermediate reasoning before a final answer improves performance on multi-step problems (Wei et al., 2022).
Asked for a final number outright, the model must commit in a single step; asked to write the intermediate quantities first, it conditions each step on the one before, which is why the gain is largest on problems with several dependent stages, such as a chained unit conversion or a conditional-probability calculation.
And giving the model \emph{tools}, letting it act, observe the result, and act again turned a text generator into a system that operates in the world, formalized as the reason-and-act loop the course builds in Chapter~\ref{chapter:AgentLoop} (Yao et al., 2023).
At that point the model is no longer only answering; it is perceiving, reasoning, and acting.

\begin{example}
A \emph{tool} is nothing more exotic than a function the surrounding program is willing to run on the model's behalf: execute this Python snippet, read this file, query this database, fetch this page.
The model itself runs nothing.
It emits a structured request; the program hosting it carries that request out and pastes the result back into the conversation as ordinary text; the model reads the result and decides what to do next.
So a question such as \emph{does the fit improve if the outliers are dropped?} stops being answered from memory.
The model writes the fit, asks for it to be run, reads the residuals that come back, notices that one predictor has gone unstable, and revises.
That loop is the entire distance between a system that describes an analysis and one that performs it.
\end{example}

The techniques for making a model reason and act well are a research field in their own right; it is worth naming the landscape even though the details belong to later chapters.
On the \emph{prompting} side, chain-of-thought has search-based extensions: \emph{self-consistency} samples several reasoning paths and keeps the majority answer, and \emph{tree-of-thoughts} searches over intermediate steps rather than committing to one.
Whereas \emph{ReAct} and \emph{Reflexion} interleave reasoning with tool calls and let a model revise after seeing a result (developed in Chapters~\ref{chapter:ContextEngineering} and~\ref{chapter:AgentLoop}).
On the \emph{training} side, reinforcement learning from human feedback (RLHF) is the instruction-tuning method named above; its relatives (learning from AI feedback, constitutional methods) and the reinforcement-learning training behind current reasoning models push the same lever further.
The shape of RLHF is worth carrying even without the details: people rank competing model responses to the same prompt, those rankings train a separate model that scores responses, and the language model is then optimized to score well, which lets a preference nobody can write down as a loss function be learned from comparisons instead.
The point is that ``make the model reason and act better'' is an active, crowded field rather than a single trick; the course develops the parts it uses and names the rest so the reader can place them.

The progression did not stop at tool use.
A recent step folds reasoning back inside the model: \emph{reasoning models}\index{Reasoning model} are trained to carry out extended chains of thought internally before they answer, rather than being prompted to reason step by step, and models are increasingly \emph{multimodal}, taking images and other data alongside text.
The practical consequences, including a latent reasoning trace that cannot be read as an audit log and a thinking budget in place of a prompt scaffold, are treated in Chapter~\ref{chapter:AgentLoop}.
Both consequences are concrete rather than abstract: the internal deliberation is optimized to help the model reach an answer, not to be read by a person, so it is a poor record of why the model did what it did; and what replaces the prompt scaffold is a coarser control, a dial setting how much deliberation to spend before answering.

Each step so far has been described from the inside, as one more limitation removed.
From outside the research community, the same sequence registered as two events.
In November~2022, the release of ChatGPT put an instruction-tuned model in front of anyone with a browser and sent a shock wave well beyond the research community, showcasing the language model as a general-purpose advisor.
Three years later, in December~2025, agentic coding crossed a practical threshold: on narrow, well-specified programming tasks whose success can be checked automatically, a frontier model such as Anthropic's Opus~4.6 approached the reliability of a junior software engineer or data scientist, turning much routine low-level coding into something a researcher directs rather than writes by hand, freeing attention for the parts of the work that call for judgment.
The qualifier carries weight: reliability falls off as tasks widen, as the specification thins, or as success stops being cheaply verifiable, so the threshold is best read as one crossed within a regime rather than across programming as a whole.
\emph{Cheaply verifiable} is the load-bearing phrase, and it means only this: some check settles the matter without a person reading the work (a test suite that passes, a simulation that reproduces a known result, a residual that drops), so an agent can iterate against it unattended, whereas without such a check someone must inspect every attempt.
Neither was a new idea so much as a new level of a capability the sequence above had been building toward.
The progression is therefore still under way, and its shape is unchanged: each step extends what the system can reliably do on its own.
Still, a legible, incremental progression is often felt from outside as a sudden arrival.

\section{The Progression Was Not Random}
\label{section:NotRandom}

Read as a list, the milestones above are a handful of loosely related papers.
Read as a chain, each one removed a specific limitation of the last.
This is a \emph{selected} teaching lineage, not a claim that each advance was the sole cause of the next; larger datasets, faster hardware, and reinforcement learning from human feedback run alongside it.
Word vectors gave meaning a computable form but handled only isolated words; sequence-to-sequence mapped whole sentences but squeezed them through one vector; the Transformer removed that bottleneck and made scale practical; pre-training made a single model's knowledge transferable across tasks; scale made that model steerable by plain text; retrieval grounded it in external, updatable knowledge; and instruction tuning and tool use let it act.
Table~\ref{table:AgenticTimeline} lists the sequence; Figure~\ref{figure:AgenticChain} draws it as the dependency chain the prose describes, each arrow labelled with the limitation the next step removed.

\begin{figure}[t]
\centering
%% agentic-chain.tex: the lineage as a causal chain (Chapter 1).
%% Milestones down the centre with their years to the left and, to the right, the
%% limitation each step removed; stacked cards on the left stand for the prehistory
%% the chapter does not retrace, faded nodes on the right for the landscape now forming.
\begin{tikzpicture}[
  font=\footnotesize, line width=1pt, >=stealth',
    box/.style={draw, rounded corners=3pt, minimum width=2.4cm, minimum height=0.8cm,
                align=center, font=\small, fill=gray!12},
    hbox/.style={draw, rounded corners=3pt, minimum width=2.4cm, minimum height=0.8cm,
                align=center, font=\small, fill=white},
    lim/.style={font=\scriptsize, text=black!55, align=left},
    yr/.style={font=\scriptsize, text=black!55},
]
\node[hbox] (h1) at (1.5,4.125) {};
\node[hbox] (h2) at (1.25,3.625) {};
\node[hbox] (h3) at (1,3.125) {};
\node[hbox] (h4) at (0.75,2.625) {};
\node[hbox] (h5) at (0.5,2.125) {};
\node[lim] at (1.5,5) {History};
\node[box] (n1) at (6,6.25) {word2vec};
\node[box] (n2) at (6,5) {seq2seq};
\node[box] (n3) at (6,3.75) {transformer};
\node[box] (n4) at (6,2.5) {pre-training};
\node[box] (n5) at (6,1.25) {GPT-3};
\node[box] (n6) at (6,0) {RAG};
\draw[->] (n1) -- (n2);
\draw[->] (n2) -- (n3);
\draw[->] (n3) -- (n4);
\draw[->] (n4) -- (n5);
\draw[->] (n5) -- (n6);
\node[lim,anchor=west] at (7.4,5.625) {isolated\\words};
\node[lim,anchor=west] at (7.4,4.375) {fixed-vector\\bottleneck};
\node[lim,anchor=west] at (7.4,3.125) {no transfer\\across tasks};
\node[lim,anchor=west] at (7.4,1.875) {task-specific\\training};
\node[lim,anchor=west] at (7.4,0.625) {frozen\\knowledge};
\node[yr,anchor=east] at (4.6,6.25) {2013};
\node[yr,anchor=east] at (4.6,5) {2014};
\node[yr,anchor=east] at (4.6,3.75) {2017};
\node[yr,anchor=east] at (4.6,2.5) {2018};
\node[yr,anchor=east] at (4.6,1.25) {2020};
\node[yr,anchor=east] at (4.6,0) {2020};
\node[yr] at (11,6.25) {2022+};
\node[box, fill=blue!15] (a0) at (11,5.375) {tuning};
\node[box, fill=blue!15] (a1) at (11.5,4.125) {CoT};
\node[box, fill=blue!15] (a2) at (12,2.875) {ReAct};
\node (f6) at (11.75,1.4375) {$\vdots$};
\node[box, fill=green!3] (a3) at (11,0) {};
\node[box, fill=green!6] (a4) at (11.25,0) {};
\node[box, fill=green!9] (a5) at (11.5,0) {};
\node[box, fill=green!12] (a6) at (11.75,0) {};
\node[box, fill=green!15] (a7) at (12,0) {Agents};
\end{tikzpicture}%

\caption{The lineage as a causal chain.
Each box is a milestone, with its year to the left; each arrow is labelled with the limitation of the step on its left that the step on its right removed.
The point is not the list of names but the dependency between them: representing isolated words made mapping whole sequences the next problem, whose fixed-vector bottleneck made attention the next, and so on until models that can act, agents.
The pale stacked cards on the left (labelled \emph{History}) stand for the much longer prehistory this chapter does not retrace; the faded nodes on the right, for the fast-growing landscape now forming around agents (new techniques, stronger pipelines, and applications across fields) that these notes sample rather than survey.}
\label{figure:AgenticChain}
\end{figure}

\begin{table}[t]
\centering
\setlength{\belowcaptionskip}{8pt}
\caption{A high-level timeline of the milestones behind agentic language models.
Each step removed a concrete limitation of the previous one; the course develops the ideas in the right-hand column in the chapters they anchor.
Years mark when each result first appeared as a preprint; the publication venues given in the Further Reading sometimes lag by a year, as with BERT (2018 preprint, NAACL~2019).}
\label{table:AgenticTimeline}
\small
\begin{tabular}{@{}l l p{6.8cm}@{}}
\toprule
Year & Breakthrough & \raggedright What it unlocked \tabularnewline
\midrule
2013  & Word embeddings                   & \raggedright meaning as a learned vector geometry \tabularnewline
2014  & Sequence-to-sequence              & \raggedright end-to-end mapping of whole sequences \tabularnewline
2017  & The Transformer                   & \raggedright attention replaces recurrence; scale becomes practical \tabularnewline
2018  & Pre-training                      & \raggedright transferable, contextual representations \tabularnewline
2020  & Scale and few-shot                & \raggedright steering a frozen model by text alone \tabularnewline
2020  & Retrieval-augmented generation    & \raggedright grounding in external, updatable knowledge \tabularnewline
2022--23 & Instruction tuning, reasoning, tool use & \raggedright following instructions and acting: agents \tabularnewline
2024--26 & Reasoning \& multimodality     & \raggedright trained latent reasoning; the chain continues \tabularnewline
\bottomrule
\end{tabular}
\end{table}

None of these steps was guaranteed in advance, and their nearness in time is why the field's current state is genuinely recent.
But the sequence is not arbitrary: today's agents are where this particular chain of problems and solutions has led, and the same pressure that drove it --- extend what the system can reliably do --- is what drives the rest of these notes.

For a researcher, this history is not a spectator's story.
On tasks that are well-specified and verifiable (writing and debugging analysis code, sweeping a method across many datasets, drafting a derivation, screening a design space), a capable agent can compress work that once took hours or days into minutes: an order-of-magnitude acceleration of the research loop, sometimes two, enough on some tasks to change what a small group can investigate independently.
The gains are uneven and should not be overstated; they shrink where judgment, genuine novelty, or hard-won trust is the bottleneck rather than throughput.
Where they do apply, though, the acceleration of the research loop is real.
It is not confined to computer science; a biologist fitting population models, an economist searching model specifications, a physicist proposing and testing functional forms, and a chemist screening candidates all face the same well-specified, checkable inner loops.

The reason to attend to this moment is therefore not novelty but \emph{discovery}: the parts of research that reduce to propose, run, check, and repeat are exactly the parts an agent can turn, and turning them faster changes the pace of finding things out.
Learning to direct that process and, by the capstone, to let agents search on their own is what these notes are for.


\section*{Further Reading}

\begin{small}
\begin{enumerate}
\item Mikolov, T., Chen, K., Corrado, G., and Dean, J., ``Efficient Estimation of Word Representations in Vector Space,'' \emph{ICLR Workshop}, 2013 (arXiv:1301.3781) — the word2vec architectures (CBOW and Skip-gram).
\item Mikolov, T., Sutskever, I., Chen, K., Corrado, G., and Dean, J., ``Distributed Representations of Words and Phrases and their Compositionality,'' \emph{NeurIPS}, 2013 (arXiv:1310.4546) — negative sampling, phrase vectors, and additive compositionality.
\item Sutskever, I., Vinyals, O., and Le, Q. V., ``Sequence to Sequence Learning with Neural Networks,'' \emph{NeurIPS}, 2014 (arXiv:1409.3215) — the recurrent encoder--decoder.
\item Bahdanau, D., Cho, K., and Bengio, Y., ``Neural Machine Translation by Jointly Learning to Align and Translate,'' \emph{ICLR}, 2015 (arXiv:1409.0473) — additive attention within a recurrent model, the first relief of the fixed-vector bottleneck.
\item Vaswani, A., et al., ``Attention Is All You Need,'' \emph{NeurIPS}, 2017 (arXiv:1706.03762) — the Transformer; attention alone, recurrence removed; the architecture underlying every model the course uses.
\item Devlin, J., Chang, M.-W., Lee, K., and Toutanova, K., ``BERT: Pre-training of Deep Bidirectional Transformers for Language Understanding,'' \emph{NAACL}, 2019 (arXiv:1810.04805) — pretraining and contextual representations.
\item Brown, T. B., et al., ``Language Models are Few-Shot Learners,'' \emph{NeurIPS}, 2020 (arXiv:2005.14165) — GPT-3 and in-context learning.
\item Lewis, P., et al., ``Retrieval-Augmented Generation for Knowledge-Intensive NLP Tasks,'' \emph{NeurIPS}, 2020 (arXiv:2005.11401) — parametric plus non-parametric memory.
\item Ouyang, L., et al., ``Training Language Models to Follow Instructions with Human Feedback,'' \emph{NeurIPS}, 2022 (arXiv:2203.02155) — instruction tuning with human feedback.
\item Wei, J., et al., ``Chain-of-Thought Prompting Elicits Reasoning in Large Language Models,'' \emph{NeurIPS}, 2022 (arXiv:2201.11903) — intermediate reasoning before an answer.
\item Yao, S., et al., ``ReAct: Synergizing Reasoning and Acting in Language Models,'' \emph{ICLR}, 2023 (arXiv:2210.03629) — interleaving reasoning with tool use.
\end{enumerate}
\end{small}
% EVOLVE-BLOCK-END
